%% kapa: exam=1 round=26 year=2015 period=1 form=A subject=civil_law subject_ko=민법 seq=2 number=2
%% source: Q-Net 감정평가사 기출문제 (cst00309) artlSeq=5138001 file=extracted/1008/5138001/제26회_감정평가사 1차_기출/제26회_감정평가사 1차_기출/1교시_감평_A형.pdf
%% conversion: kordoc 4.24.51 -> tools/gen1.py (auto); review=proofread-sonnet; flags=none
\begin{question}{2}
신의칙에 관한 설명으로 옳지 않은 것은? (다툼이 있으면 판례에 따름)
\begin{choices}
  \item 신의칙 위반 여부는 당사자의 주장이 없더라도 법원이 직권으로 판단할 수 있다.
  \item 부동산 점유자가 취득시효완성 후에 그 사실을 모르고 소유자에게 당해 토지에 관하여 어떠한 권리 도 주장하지 않기로 하였는데, 나중에 시효완성을 주장하는 것은 특별한 사정이 없는 한 신의칙에 반한다.
  \item 법률행위가 법령에 위반되어 무효임을 알면서도 그 법률행위를 한 자가 나중에 강행법규 위반을 이유로 무효를 주장하더라도 신의칙에 반하지 않는다.
  \item 매매계약체결 후 9년이 지났고 시가가 올랐다는 사정만으로 매수인의 소유권 이전등기절차 이행청구가 신의칙에 위배된다고 할 수 없다.
  \item 부동산 거래에 있어 신의칙상 상대방에게 고지의무의 대상이 되는 것은 법령의 규정뿐이고, 널리 계약상, 관습상 또는 조리상의 일반원칙에 의해서는 인정될 수 없다.
\end{choices}
\end{question}
\answer{5}
